\section{Experiment 14B: Ten-Class Fashion-MNIST and Multiclass Event Learning}
\label{sec:exp14b}

Experiment~14A established that sparse temporal gradient events can train a $25665$-parameter binary Fashion-MNIST classifier under asynchronous label-skew heterogeneity. Experiment~14B increases the task complexity from one binary separator to the complete ten-class Fashion-MNIST problem~\cite{xiao2017fashionmnist}. The purpose is to test whether the same event mechanism survives multiclass softmax competition and class-dependent client objectives before introducing convolutional architectures, clustering, personalization, or a new communication rule.

A deliberate design choice is to keep the hidden architecture almost unchanged from Experiment~14A. Rather than simultaneously increasing hidden width and the number of classes, we use
\begin{equation}
784\rightarrow32\;\tanh\rightarrow16\;\tanh\rightarrow10,
\end{equation}
with
\begin{equation}
\boxed{d=25818}
\end{equation}
trainable parameters. Hence the principal experimental change is the binary-to-multiclass transition rather than another large increase in parameter dimension. This isolates the effect of ten coupled output logits and ten-way label heterogeneity.

\subsection{Questions and scientific gates}

Experiment~14B asks five questions.
\begin{enumerate}
    \item Can raw-threshold LIF gradient communication optimize a ten-class softmax network without a class-specific or layer-specific communication mechanism?
    \item At useful multiclass accuracy, does the event representation remain substantially cheaper than dense float32 communication and competitive with EF-TopK?
    \item Does stronger label skew produce a permanent multiclass failure, or mainly change the communication demand and convergence transient?
    \item Are aggregate accuracy and cross-entropy sufficient to characterize success, or does the sparse event process create class-wise convergence failures that are hidden by the mean metric?
    \item Does coupling a class-dominant client to a particular asynchronous compute period create a new fairness sensitivity despite the existing $p_iT_i$ rate normalization?
\end{enumerate}

The experiment is explicitly not a clustering test. Every client still participates in one common globally balanced ten-class objective. Clustering is justified only if a genuinely multimodal or incompatible-client failure appears, not merely because individual clients have different label frequencies.

\subsection{Exact multiclass federation}

Fashion-MNIST contains $6000$ training examples for each of ten classes. All $60000$ official training images and all $10000$ official test images are used. Each of the ten clients receives exactly $6000$ training images, and each class contributes exactly $6000$ images globally. This lets label skew vary without changing client sample count or global class balance.

The client completion periods remain
\begin{equation}
T=[1,1,2,2,5,5,10,10,20,20],
\end{equation}
with objective weights $p_i=0.1$. In the default construction, client $i$ is dominant in Fashion-MNIST class $i$. Four label-skew regimes are defined:
\begin{itemize}
    \item \textbf{IID:} $600$ examples of every class per client;
    \item \textbf{moderate:} $2400$ examples of one dominant class and $400$ of each other class, so the dominant fraction is $40\%$;
    \item \textbf{strong:} $3300$ dominant-class examples and $300$ of each other class, so the dominant fraction is $55\%$;
    \item \textbf{extreme:} $5100$ dominant-class examples and $100$ of each other class, so the dominant fraction is $85\%$.
\end{itemize}
Each count matrix has row and column sums equal to $6000$. Therefore the global empirical class distribution remains exactly uniform in every regime.

Partition seeds change the concrete images assigned to clients while preserving the count matrix. A separate permutation audit changes which semantic class is dominant on each compute period. This isolates label-skew severity from the accidental pairing between a particular Fashion-MNIST class and a fast or slow client.

\subsection{Multiclass objective and gradient}

For sample $(x,y)$ with $y\in\{0,\ldots,9\}$, the MLP produces logits
\begin{equation}
\ell(x;w)\in\mathbb{R}^{10}.
\end{equation}
The class probabilities are
\begin{equation}
\pi_c(x;w)=
\frac{\exp(\ell_c(x;w))}
{\sum_{r=0}^{9}\exp(\ell_r(x;w))},
\end{equation}
and the predictive loss is the standard softmax cross-entropy
\begin{equation}
\mathcal{L}_{\mathrm{CE}}(w)
=-\frac{1}{B}\sum_{b=1}^{B}
\log \pi_{y_b}(x_b;w).
\end{equation}
The local stochastic objective includes the same $\ell_2$ regularization as Experiment~14A,
\begin{equation}
F_i(w)=\mathcal{L}_{\mathrm{CE},i}(w)+\frac{\lambda}{2}\norm{w}^2,
\qquad \lambda=10^{-4},
\end{equation}
with mini-batch size $B=32$. Hidden layers use $\tanh$ activations. All reported model-selection sweeps use training objective values rather than held-out test accuracy.

In addition to test cross-entropy and overall accuracy, Experiment~14B reports the worst class accuracy
\begin{equation}
A_{\min}=\min_{c\in\{0,\ldots,9\}} A_c.
\end{equation}
Because the official test set is exactly balanced across classes, macro accuracy equals overall accuracy. The minimum class accuracy is nevertheless essential: the early audits show that a model can exceed $70\%$ mean accuracy while one class is still almost completely unlearned.

\subsection{Communication rules and accounting}

The event rule is unchanged in structure:
\begin{equation}
z_i^{k+1}=\rho z_i^k-\gamma p_iT_i g_i^k,
\end{equation}
followed by a full reset of coordinates satisfying
\begin{equation}
|z_{i,j}^{k+1}|\ge\Delta.
\end{equation}
For a fired coordinate,
\begin{equation}
w_j^+=w_j+q(t)\sign(z_{i,j}),
\end{equation}
with
\begin{equation}
q(t)=q_0\left(1+\frac{t}{500}\right)^{-0.1}.
\end{equation}
No descent certificate, local personalization, client clustering, class reweighting, layer-wise threshold, or learned compressor is used.

Since $d=25818$,
\begin{equation}
B_{\mathrm{event}}
=\lceil\log_2 25818\rceil+1
=16\ \text{bits}
\end{equation}
are charged for every coordinate event: $15$ address bits and one sign bit. A dense float32 vector requires
\begin{equation}
B_{\mathrm{dense}}=32d=826176\ \text{bits}.
\end{equation}
EF-TopK sends a float32 residual value and address for every selected coordinate. As before, the $32$-bit illustrative packet header is tracked separately from the logical payload. The comparisons below use logical uplink payload unless stated otherwise.

\subsection{Why the 250-tick audit was not a valid endpoint}

A first 250-tick strong-skew audit verified that all three optimizer families were trainable but also exposed a multiclass validity problem. Dense SGD with the best tested short-audit step reached $69.99\%$ overall accuracy, yet its worst-class accuracy was only $27.8\%$. EF-TopK and the initial event points were more class-starved, with several configurations still near zero worst-class accuracy. The best initial event gain, $\gamma=0.4$, reached $58.73\%$ mean accuracy but still had an unlearned class.

This means that a short multiclass run can give a superficially reasonable average accuracy while hiding a severe class failure. Experiment~14B therefore does not use the 250-tick result as a performance comparison. The horizon was extended and class-wise accuracy became a formal gate.

\subsection{650-tick optimizer and event-scale refinement}

At $650$ ticks, dense and EF-TopK step sizes were refined. Dense SGD with step $0.08$ gives the lowest final training objective among $\{0.02,0.04,0.08\}$ at this horizon. EF-TopK with fractions $0.5\%$ and $2.5\%$ is also best by final training objective at step $0.08$ among the tested $\{0.04,0.08\}$ values.

The first event refinement showed a monotone benefit from increasing evidence gain beyond the binary Experiment~14A value. A focused sweep then used
\begin{equation}
\gamma\in\{0.6,0.8,1.0\},
\qquad
q_0\in\{0.0035,0.005,0.0075\},
\qquad
\Delta=0.025.
\end{equation}
The selected configuration is
\begin{equation}
\boxed{
\rho=0.999,\qquad
\gamma=1.0,\qquad
\Delta=0.025,\qquad
q_0=0.0035.
}
\end{equation}
It has the lowest final training objective in this focused event sweep. At the same $\gamma=1$, increasing the jump to $q_0=0.0075$ degrades test cross-entropy to approximately $0.883$ and accuracy to $64.45\%$. Thus the positive result is not obtained by simply maximizing firing rate or event amplitude. Evidence accumulation and jump resolution remain coupled optimization parameters.

\subsection{Primary strong-skew result at 650 ticks}

For the default strong partition with seed $2400$, the selected event model reaches the following representative operating point.
\begin{takeawaybox}[title={Experiment 14B primary event result}]
Test CE is $0.64007$, overall accuracy is $78.23\%$, worst-class accuracy is $35.3\%$, and logical uplink payload is $34.75$~Mbit.
\end{takeawaybox}
The matched references are summarized in Table~\ref{tab:exp14b-primary}.

\begin{table}[t]
\centering
\caption{Experiment~14B at $650$ ticks on the default strong label-skew partition. Baseline steps are selected from the stated training-objective sweeps.}
\label{tab:exp14b-primary}
\begin{tabular}{lrrrr}
\toprule
Method & Test CE & Accuracy & Worst class & Payload\\
\midrule
EF-TopK $0.5\%$ & 0.74371 & 0.7219 & 0.198 & 14.58 M\\
EF-TopK $2.5\%$ & 0.69286 & 0.7246 & \textbf{0.378} & 72.88 M\\
Dense SGD & 0.67331 & 0.7496 & 0.355 & 1986.13 M\\
\textbf{Events} & \textbf{0.64007} & \textbf{0.7823} & 0.353 & \textbf{34.75 M}\\
\bottomrule
\end{tabular}
\end{table}

On this partition, event communication uses about
\begin{equation}
57.2\times
\end{equation}
fewer bits than dense SGD and about
\begin{equation}
2.10\times
\end{equation}
fewer bits than EF-TopK $2.5\%$. It also has lower cross-entropy and higher mean accuracy than both references. Its worst-class accuracy is essentially equal to dense SGD and $2.5$ percentage points below EF-TopK $2.5\%$.

This single partition is a strong result, but the independent-partition audit below prevents it from being promoted to a uniform Pareto-dominance claim.

\subsection{Independent strong-partition robustness}

The selected event hyperparameters are frozen and repeated for partition seeds $2400$, $2500$, and $2600$. Table~\ref{tab:exp14b-seeds} reports the observed results.

\begin{table}[t]
\centering
\caption{Experiment~14B robustness across independent strong label-skew partitions.}
\label{tab:exp14b-seeds}
\begin{tabular}{llrrr}
\toprule
Seed & Method & Test CE & Accuracy & Worst class\\
\midrule
2400 & Dense & 0.6733 & 0.7496 & 0.355\\
     & EF $2.5\%$ & 0.6929 & 0.7246 & 0.378\\
     & Events & \textbf{0.6401} & \textbf{0.7823} & 0.353\\
\midrule
2500 & Dense & \textbf{0.7339} & 0.7175 & \textbf{0.343}\\
     & EF $2.5\%$ & 0.7375 & \textbf{0.7271} & 0.331\\
     & Events & 0.7915 & 0.6977 & 0.082\\
\midrule
2600 & Dense & \textbf{0.6265} & \textbf{0.7706} & 0.345\\
     & EF $2.5\%$ & 0.6454 & 0.7492 & 0.174\\
     & Events & 0.6593 & 0.7654 & \textbf{0.361}\\
\bottomrule
\end{tabular}
\end{table}

The event means over these three partitions are
\begin{equation}
\text{CE}=0.6970\pm0.0825,
\qquad
A=0.7485\pm0.0448,
\qquad
A_{\min}=0.265\pm0.159,
\end{equation}
with approximately $34.3$ Mbit mean payload. EF-TopK $2.5\%$ has mean CE $0.6919\pm0.0461$, accuracy $0.7336\pm0.0135$, and worst-class accuracy $0.294\pm0.107$ at $72.88$ Mbit. Dense SGD has mean CE $0.6779\pm0.0538$, accuracy $0.7459\pm0.0267$, and a much more stable worst-class accuracy $0.348\pm0.006$ at $1.986$ Gbit.

This is a central negative result. The event method remains communication-efficient and its mean multiclass performance is competitive, but the worst-class convergence is more variable. Seed $2500$ is an explicit failure at the fixed $650$-tick horizon: overall accuracy is still about $69.8\%$, while the weakest class reaches only $8.2\%$. Therefore
\begin{equation}
\boxed{\text{uniform finite-horizon class-wise robustness: NOT ESTABLISHED}.}
\end{equation}
Aggregate accuracy alone would hide this failure.

\subsection{Long-horizon strong-skew audit}

A $1200$-tick audit distinguishes persistent class failure from slow convergence. It also reveals that the constant baseline step selected at $650$ ticks is not automatically suitable at the longer horizon. Dense step $0.08$, for example, degrades by $1200$ ticks, whereas steps $0.02$--$0.04$ remain stable. The long-horizon result must therefore be compared against a separate baseline step audit rather than mechanically reusing the short-horizon optimum.

Table~\ref{tab:exp14b-long} summarizes the important long-horizon points on seed $2400$.

\begin{table}[t]
\centering
\caption{Experiment~14B strong-skew result at $1200$ ticks.}
\label{tab:exp14b-long}
\begin{tabular}{lrrrr}
\toprule
Method & Test CE & Accuracy & Worst class & Payload\\
\midrule
Dense, step $0.02$ & 0.70803 & 0.7462 & 0.491 & 3668.22 M\\
Dense, step $0.04$ & 0.70197 & 0.7363 & 0.440 & 3668.22 M\\
EF $0.5\%$, step $0.04$ & 0.68752 & 0.7323 & 0.305 & 26.92 M\\
EF $2.5\%$, step $0.04$ & 0.62352 & 0.7712 & 0.457 & 134.60 M\\
EF $2.5\%$, step $0.08$ & \textbf{0.61715} & 0.7685 & \textbf{0.560} & 134.60 M\\
\textbf{Events} & 0.62883 & \textbf{0.7824} & 0.416 & \textbf{66.99 M}\\
\bottomrule
\end{tabular}
\end{table}

The event model now uses about $54.8\times$ fewer payload bits than dense communication and approximately $2.01\times$ fewer than EF-TopK $2.5\%$. Its final training objective $0.6201$ is essentially the same as EF-TopK $2.5\%$ at $0.6155$, while its mean test accuracy is $1.39$ percentage points higher. EF-TopK retains a clear worst-class advantage, $56.0\%$ versus $41.6\%$, and a slightly lower test cross-entropy. Hence the long-horizon result is a communication-efficiency pass, but not a strict domination of the complete EF performance vector.

The event per-class accuracies are
\begin{equation}
[0.416,0.898,0.795,0.877,0.720,0.835,0.557,0.836,0.941,0.949].
\end{equation}
The weakest classes are class~0 (T-shirt/top) and class~6 (Shirt), which are also visually confusable Fashion-MNIST categories. This observation is descriptive rather than a proof that visual ambiguity causes the event transient.

\subsection{Label-skew severity and endogenous communication}

With the selected strong-skew event hyperparameters frozen, logical event payload increases markedly with label skew:
\begin{equation}
10.54\ \text{Mbit (IID)}
\rightarrow
25.06\ \text{Mbit (moderate)}
\rightarrow
34.75\ \text{Mbit (strong)}
\rightarrow
55.53\ \text{Mbit (extreme)}.
\end{equation}
Thus client disagreement again appears as increased endogenous event traffic rather than as a fixed communication budget.

The predictive comparison is not uniformly favorable to events. Under IID data, EF-TopK $0.5\%$ and $2.5\%$ reach test CE $0.5938$ and $0.5698$, respectively, while the frozen event rule reaches $0.6723$. Under moderate skew, events reach CE $0.6636$ and $76.1\%$ accuracy. Under strong skew, the event point becomes especially competitive. This gives an important qualification:
\begin{equation}
\boxed{\text{the event advantage is regime-dependent, not a universal sparse-optimizer dominance}.}
\end{equation}
The event parameters were selected on the strong regime, so this sweep is a transfer test rather than an IID retuning campaign.

\subsection{Extreme skew: transient failure versus persistent failure}

At the fixed $650$-tick horizon with $85\%$ dominant clients, the default event run gives $69.61\%$ mean accuracy but only $2.4\%$ worst-class accuracy. A naive interpretation would be that extreme label skew breaks the multiclass event mechanism. Two audits show that this conclusion is too strong.

First, permuting which semantic class is dominant on each compute period changes the $650$-tick event worst-class accuracy from $2.4\%$ in the identity assignment to $36.2\%$ in the reverse assignment and $0\%$ in one mixed assignment. Dense and EF-TopK also show substantial sensitivity under these extreme rotations. The effect therefore involves the interaction of label skew, class difficulty, asynchronous compute periods, and optimizer dynamics rather than a simple event-only failure.

Second, extending the default extreme split to $1200$ ticks gives the results in Table~\ref{tab:exp14b-extreme-long}.

\begin{table}[t]
\centering
\caption{Experiment~14B extreme $85\%$ label skew at $1200$ ticks.}
\label{tab:exp14b-extreme-long}
\begin{tabular}{lrrrr}
\toprule
Method & Test CE & Accuracy & Worst class & Payload\\
\midrule
Dense, step $0.04$ & \textbf{0.60256} & \textbf{0.7872} & \textbf{0.538} & 3668.22 M\\
EF $2.5\%$, step $0.08$ & 0.96294 & 0.6919 & 0.390 & 134.60 M\\
Events & 0.68202 & 0.7849 & 0.476 & \textbf{101.49 M}\\
\bottomrule
\end{tabular}
\end{table}

The event worst class recovers from $2.4\%$ to $47.6\%$, while overall accuracy rises to $78.49\%$, essentially matching dense accuracy at about $36.1\times$ lower payload. The event model also substantially outperforms the tested EF-TopK $2.5\%$ point in cross-entropy, mean accuracy, and worst-class accuracy while using fewer bits. Thus extreme skew produces a serious finite-horizon class imbalance, but it does not establish a permanent inability to learn the ten-class objective.

\subsection{Packet density and tensor activity}

Multiclass learning requires a denser event stream than the binary Experiment~14A. On the selected strong $650$-tick run, $2171991$ coordinate events are transmitted in $2404$ nonempty messages, or approximately
\begin{equation}
903.5
\end{equation}
coordinates per packet. This is about $3.50\%$ of the $25818$-parameter model. The representation therefore remains sparse, but the natural scalable communication object is now clearly a sparse event packet rather than an isolated coordinate spike.

The selected configuration reaches essentially complete parameter coverage by $650$ ticks. This contrasts with binary Experiment~14A, where only about $65\%$ of parameters ever fired. The multiclass objective creates enough diverse gradient evidence that nearly every coordinate participates at least once. Tensor-wise firing statistics are reported in the reproducible diagnostic driver and are used to check for output-layer or hidden-layer starvation. No class-specific output threshold is introduced in this experiment.

\subsection{Negative results and what they mean}

The most informative Experiment~14B negative findings are:
\begin{enumerate}
    \item \textbf{Average accuracy can hide class collapse.} At short horizons, models with apparently respectable mean accuracy may have near-zero accuracy on at least one class. Worst-class accuracy must therefore remain a primary metric.
    \item \textbf{The binary event scale does not transfer directly.} Multiclass learning requires substantially larger evidence gain. The selected $\gamma=1.0$ is more than three times the binary 14A value $0.3$.
    \item \textbf{Larger event jumps can destabilize the multiclass model.} At $\gamma=1.0$, $q_0=0.0075$ performs far worse than $q_0=0.0035$.
    \item \textbf{Strong-partition robustness is not uniform.} One independent partition has only $8.2\%$ worst-class accuracy at $650$ ticks despite reasonable aggregate performance.
    \item \textbf{The event rule does not dominate EF-TopK under every regime.} The frozen strong-skew configuration is weaker than EF on IID data.
    \item \textbf{Extreme-skew class collapse is largely transient, but the transient is important.} The weakest class can be almost absent at $650$ ticks and recover strongly by $1200$ ticks.
    \item \textbf{The best constant dense step depends on horizon.} A step that is favorable at $650$ ticks can degrade by $1200$ ticks. Baseline fairness therefore requires horizon-specific trainability checks.
\end{enumerate}

These findings shift the main scaling concern. The question is no longer whether sparse temporal events can train a multiclass real-data network. They can. The open issue is whether the event architecture provides sufficiently predictable \emph{class-wise convergence transients} under the combination of label skew and asynchronous compute.

\subsection{Experiment 14B verdict}

The resulting gates are
\begin{equation}
\boxed{\text{ten-class Fashion-MNIST event-learning viability: PASS}},
\end{equation}
\begin{equation}
\boxed{\text{strong-skew communication efficiency: PASS}},
\end{equation}
\begin{equation}
\boxed{\text{long-horizon performance at substantially reduced bits: PASS}},
\end{equation}
\begin{equation}
\boxed{\text{uniform Pareto dominance across regimes and partitions: FAIL / NOT ESTABLISHED}},
\end{equation}
\begin{equation}
\boxed{\text{finite-horizon worst-class robustness: QUALIFIED / OPEN}},
\end{equation}
\begin{equation}
\boxed{\text{permanent failure under extreme label skew: NOT OBSERVED}}.
\end{equation}

The central empirical result is therefore more nuanced than Experiment~14A:
\begin{equation}
\boxed{
\text{multiclass sparse-event learning scales, but class-wise convergence is more fragile than aggregate convergence.}
}
\end{equation}
The mechanism remains highly communication-efficient on the strong and long-horizon tests, yet a paper-quality argument must explicitly include its partition-sensitive transient behavior rather than report only mean test accuracy.

The reusable implementation is in \texttt{src/neuromorphicfl/fmnist\_multiclass\_benchmark.py}. The Experiment~14B audit, refinement, fairness, final robustness, and layer-diagnostic drivers are in \texttt{experiments/14b\_*.py}. Report-ready observed tables are stored under \texttt{experiments/results/14b\_fmnist\_multiclass/}.
